\section{Feature Selection}\label{sec:feature_selection}

Proses seleksi fitur (\textit{feature selection}) bertujuan untuk memilih bagian data (kolom) yang paling berkaitan dengan target akhir. Pemilihan ini penting untuk membuat hasil prediksi lebih akurat, mempercepat proses perhitungan komputer, dan mencegah model hanya menghafal data tanpa memahami polanya (\textit{overfitting}). Secara umum, pemilihan fitur dibagi menjadi tiga cara, yaitu berdasarkan pengetahuan dari pakar (\textit{domain knowledge}), perhitungan statistik (\textit{filter methods}), dan pendekatan heuristik atau pencarian terarah (\textit{wrapper methods}).

\subsection{Pendekatan Pengetahuan Pakar (\textit{Domain Knowledge})}
Sebelum melakukan pemilihan data kepada algoritma komputer, langkah pertama yang paling penting adalah menyaring data menggunakan pengetahuan pakar di lapangan (\textit{domain knowledge}). Cara ini dilakukan dengan menentukan informasi apa saja yang secara pengalaman nyata berkaitan erat dengan apa yang ingin ditebak, misalnya melalui wawancara mendalam dengan Pengelola apartemen atau petugas keamanan.

Menurut Kuhn dan Johnson  \cite{kuhn2013applied}, pemanfaatan pengetahuan domain sering kali lebih penting daripada teknik pemilihan fitur otomatis mana pun. Hal ini dikarenakan algoritma seleksi fitur murni secara matematis berpotensi memilih variabel yang hanya memiliki korelasi kebetulan (\textit{spurious correlation}) tanpa adanya dasar hubungan sebab-akibat yang logis. Untuk mendeteksi penyewa yang bermasalah, wawasan dari para pakar lapangan ini sangat membantu menentukan fitur khusus---seperti ada tidaknya riwayat pelanggaran atau ketidakwajaran waktu \textit{check-in}---sebagai tanda yang kuat (\textit{strong signal}) adanya aktivitas yang tidak normal.

\subsection{Metode Statistik (Filter Methods)}
Setelah fitur-fitur awal ditentukan, metode statistik digunakan untuk menghitung seberapa kuat hubungan antara fitur tersebut dengan status penyewa secara matematis. Mengingat pada kasus ini jumlah penyewa yang tidak normal (1) jauh lebih sedikit dibandingkan yang normal (0), perhitungan berikut merupakan adaptasi konseptual dari buku dasar \textit{Data Mining} \cite{han2011data} dan teori informasi \cite{cover2006elements} yang disimulasikan menggunakan perbandingan data tidak seimbang (\textit{imbalanced}).
    
\begin{enumerate}
    \item \textbf{Chi-Square Test} \\
    Uji Chi-Square ($\chi^2$) digunakan untuk melihat seberapa kuat hubungan antara dua data berbentuk kategori, seperti metode pembayaran dengan status penyewa. Menurut \cite{han2011data}, rumusnya adalah sebagai berikut:
    $$
    \chi^2 = \sum \frac{(O_i - E_i)^2}{E_i}
    $$
    dengan:
    \begin{itemize}
        \item $O_i$ : nilai observasi (angka yang sebenarnya terjadi),
        \item $E_i$ : nilai harapan (angka teoretis jika tidak ada hubungan).
    \end{itemize}

    Misalkan kita ingin menguji hubungan metode pembayaran (tunai/kartu) dengan status penyewa (normal/tidak normal) di mana kelas tidak normal merupakan minoritas.

    \begin{table}[H]
    \centering
    \caption{Tabel Kontingensi Metode Pembayaran dan Status Penyewa}
        \begin{tabular}{|c|c|c|c|}
            \hline
            \textbf{Metode Pembayaran} & \textbf{Normal (0)} & \textbf{Tidak Normal (1)} & \textbf{Total} \\
            \hline
            Tunai & 40 & 15 & 55 \\
            Kartu & 40 & 5 & 45 \\
            \hline
            Total & 80 & 20 & 100 \\
            \hline
        \end{tabular}
    \end{table}

    Langkah-langkah perhitungannya adalah sebagai berikut. Nilai observasi ($O_{ij}$) seperti angka 40, 15, dan 5 didapatkan langsung dari perhitungan data riil pada tabel kontingensi di atas. Kemudian, nilai harapan ($E_{ij}$) untuk masing-masing sel dihitung menggunakan rumus:
    $$
    E_{ij} = \frac{\text{(Total Baris)} \times \text{(Total Kolom)}}{\text{Grand Total}}
    $$
    
    Berikut adalah rincian perhitungan nilai harapan ($E_i$) untuk setiap kombinasi kategori:
    \begin{itemize}
        \item \textbf{Tunai dan Normal}:
        $$
        E_{\text{Tunai, Normal}} = \frac{55 \times 80}{100} = 44
        $$
        \item \textbf{Tunai dan Tidak Normal}:
        $$
        E_{\text{Tunai, Tidak Normal}} = \frac{55 \times 20}{100} = 11
        $$
        \item \textbf{Kartu dan Normal}:
        $$
        E_{\text{Kartu, Normal}} = \frac{45 \times 80}{100} = 36
        $$
        \item \textbf{Kartu dan Tidak Normal}:
        $$
        E_{\text{Kartu, Tidak Normal}} = \frac{45 \times 20}{100} = 9
        $$
    \end{itemize}

    Berdasarkan rincian nilai $O_i$ dan $E_i$ tersebut, perhitungan total nilai $\chi^2$ adalah sebagai berikut.
    $$
    \chi^2 = \frac{(40 - 44)^2}{44} + \frac{(15 - 11)^2}{11} + \frac{(40 - 36)^2}{36} + \frac{(5 - 9)^2}{9}
    $$
    $$
    \chi^2 \approx 0{,}36 + 1{,}45 + 0{,}44 + 1{,}77 = 4{,}02
    $$

    Semakin besar angka $\chi^2$, semakin kuat hubungan antara fitur tersebut dengan target hasil. Nilai $p$ dari distribusi hasil ini kemudian bisa digunakan untuk memastikan apakah hubungan tersebut benar-benar dapat dipercaya secara statistik (signifikan).

    \item \textbf{Korelasi Pearson} \\
    Korelasi Pearson digunakan untuk mengukur hubungan berupa garis lurus (linier) antara data berbentuk angka (misalnya lama menginap) dengan status penyewa \cite{han2011data}. Berikut rumusnya.
    $$
    r = \frac{\sum (x_i - \bar{x})(y_i - \bar{y})}{\sqrt{\sum (x_i - \bar{x})^2 \sum (y_i - \bar{y})^2}}, \quad r \in [-1,1]
    $$
    
    Sebagai contoh, pada sampel acak berikut penyewa tidak normal (1) muncul lebih sedikit:
    \begin{table}[H]
    \centering
    \caption{Data Lama Menginap dan Status}
    \begin{tabular}{|c|c|}
    \hline
    \textbf{Lama Menginap} ($x$) & \textbf{Status Tidak Normal} ($y$) \\
    \hline
    2 & 0 \\
    3 & 0 \\
    5 & 0 \\
    1 & 1 \\
    10 & 0 \\
    \hline
    \end{tabular}
    \end{table}
    
    \begin{enumerate}
        \item Hitung rata-rata untuk $x$ dan $y$.
        $$
        \bar{x} = \frac{2 + 3 + 5 + 1 + 10}{5} = \frac{21}{5} = 4{,}2, \quad
        \bar{y} = \frac{0 + 0 + 0 + 1 + 0}{5} = \frac{1}{5} = 0{,}2
        $$
    
        \item Hitung selisih nilai dari rata-rata, kuadratnya, dan hasil kali.
        \begin{table}[H]
        \centering
        \caption{Perhitungan Komponen Korelasi Pearson}
        \begin{tabular}{|c|c|c|c|c|c|c|}
        \hline
        $x_i$ & $y_i$ & $x_i - \bar{x}$ & $y_i - \bar{y}$ & $(x_i - \bar{x})^2$ & $(y_i - \bar{y})^2$ & $(x_i - \bar{x})(y_i - \bar{y})$ \\
        \hline
        2 & 0 & -2{,}2 & -0{,}2 & 4{,}84 & 0{,}04 & 0{,}44 \\
        3 & 0 & -1{,}2 & -0{,}2 & 1{,}44 & 0{,}04 & 0{,}24 \\
        5 & 0 & 0{,}8 & -0{,}2 & 0{,}64 & 0{,}04 & -0{,}16 \\
        1 & 1 & -3{,}2 & 0{,}8 & 10{,}24 & 0{,}64 & -2{,}56 \\
        10 & 0 & 5{,}8 & -0{,}2 & 33{,}64 & 0{,}04 & -1{,}16 \\
        \hline
        \multicolumn{4}{|c|}{\textbf{Total}} & 50{,}80 & 0{,}80 & -3{,}20 \\
        \hline
        \end{tabular}
        \end{table}
    
        \item Hitung nilai Pearson $r$.
        $$
        r = \frac{-3{,}20}{\sqrt{50{,}80 \times 0{,}80}} = \frac{-3{,}20}{\sqrt{40{,}64}} \approx \frac{-3{,}20}{6{,}37} \approx -0{,}502
        $$
    \end{enumerate}
    
    Nilai $r = -0{,}502$ menunjukkan adanya hubungan negatif berskala sedang antara lama menginap dan status pada kelompok data ini. Artinya, semakin singkat seseorang menginap, probabilitas ia dicurigai sebagai penyewa tidak normal cenderung semakin tinggi.
    
    \item \textbf{Mutual Information (MI)} \\
    Berdasarkan teori informasi \cite{cover2006elements}, MI mengukur seberapa banyak informasi (petunjuk) yang bisa didapatkan tentang target akhir jika kita mengetahui sebuah fitur tertentu, baik hubungannya berupa garis lurus maupun tidak. Metode ini cocok untuk data angka yang sudah dikelompokkan. Berikut rumusnya:
    $$
    I(X;Y) = \sum_{x \in X} \sum_{y \in Y} p(x,y) \cdot \log_2 \left( \frac{p(x,y)}{p(x)p(y)} \right)
    $$
    dengan:
    \begin{itemize}
        \item $p(x,y)$ : peluang kemunculan bersama antara fitur dan target,
        \item $p(x), p(y)$ : peluang kemunculan masing-masing secara terpisah.
    \end{itemize}
    
    Misalkan fitur lama menginap dikelompokkan menjadi: $\le$ 2 hari, 3--5 hari, dan $>$ 5 hari. Distribusi populasinya didominasi oleh kelas normal (0):
    
    \begin{table}[H]
    \centering
    \caption{Frekuensi Gabungan Lama Menginap dan Status Penyewa}
    \begin{tabular}{|c|c|c|}
    \hline
    \textbf{Lama menginap} & \textbf{Status} & \textbf{Frekuensi} \\
    \hline
    $\leq$ 2 hari & Normal (0) & 25 \\
    $\leq$ 2 hari & Tidak Normal (1) & 15 \\
    3--5 hari & Normal (0) & 35 \\
    3--5 hari & Tidak Normal (1) & 5 \\
    $>$ 5 hari & Normal (0) & 15 \\
    $>$ 5 hari & Tidak Normal (1) & 5 \\
    \hline
    \multicolumn{2}{|c|}{\textbf{Total}} & 100 \\
    \hline
    \end{tabular}
    \end{table}
    
    \begin{enumerate}
        \item Hitung probabilitas gabungan $p(x, y)$:
        $$
        \begin{aligned}
        p(\leq2, \text{Normal}) &= 0{,}25, \quad &p(\leq2, \text{Tidak normal}) &= 0{,}15 \\
        p(3{-}5, \text{Normal}) &= 0{,}35, \quad &p(3{-}5, \text{Tidak normal}) &= 0{,}05 \\
        p(>5, \text{Normal}) &= 0{,}15, \quad &p(>5, \text{Tidak normal}) &= 0{,}05
        \end{aligned}
        $$
    
        \item Hitung probabilitas individual:
        $$
        \begin{aligned}
        p(\leq2) &= 0{,}40, \quad p(3{-}5) = 0{,}40, \quad p(>5) = 0{,}20 \\
        p(\text{Normal}) &= 0{,}75, \quad p(\text{Tidak normal}) = 0{,}25
        \end{aligned}
        $$
    
        \item Hitung total nilai MI dari peluang di atas:
        $$
        \begin{aligned}
        MI = & \left(0{,}25 \cdot \log_2 \frac{0{,}25}{0{,}40 \cdot 0{,}75}\right) + \left(0{,}15 \cdot \log_2 \frac{0{,}15}{0{,}40 \cdot 0{,}25}\right) + \\
             & \left(0{,}35 \cdot \log_2 \frac{0{,}35}{0{,}40 \cdot 0{,}75}\right) + \left(0{,}05 \cdot \log_2 \frac{0{,}05}{0{,}40 \cdot 0{,}25}\right) + \\
             & \left(0{,}15 \cdot \log_2 \frac{0{,}15}{0{,}20 \cdot 0{,}75}\right) + \left(0{,}05 \cdot \log_2 \frac{0{,}05}{0{,}20 \cdot 0{,}25}\right)
        \end{aligned}
        $$
        $$
        MI \approx -0{,}065 + 0{,}087 + 0{,}077 - 0{,}050 + 0 + 0 \approx 0{,}049
        $$
    \end{enumerate}
    
    Nilai MI sebesar 0,049 menunjukkan bahwa fitur lama menginap hanya memberikan sedikit informasi tambahan untuk menebak status penyewa pada kelompok data yang sangat timpang ini.

    
